% Part: proof-theory
% Chapter: normalization
% Section: reductions

\documentclass[../../../include/open-logic-section]{subfiles}

\begin{document}

\olfileid{pt}{nor}{red}

\olsection{Reduction Conversions}

If a cut in !!a{proof} in \Log{N1i} has length~$1$, it is a segment
consisting of a single !!{formula} occurrence, and that !!{formula} is
both the conclusion of !!a{introduction} inference and the major
premise of !!a{elimination} inference. In the normalization proof,
such cuts are ``reduced'' by replacing the sub-!!{proof} ending in
such an !!{introduction}/!!{elimination} pair by a new !!{proof} in
which the pair of inferences has been eliminated. Such a replacement
may result in new cuts. However, we can convince ourselves that the
new cuts have lower rank than the cut being reduced. Thus, the new
!!{proof} either has lower cut length (because a maximal cut of
length~$1$ has been removed) or the cut rank of the !!{proof} has been
reduced (if the maximal cut in question was the only maximal cut).

For instance, if the inferences involved are \Intro{\lif} and
\Elim{\lif}, $!A \ident !B \lif !C$ and there is a
sub-!!{proof}~$\delta_1$ that ends in:
\[
\AxiomC{$\Discharge{!B}{x}$}
\RightLabel{$\delta_2$}
\DeduceC{$!C$}
\DischargeRule{\Intro{\lif}}{x}
\UnaryInfC{$!B \lif !C$}
\AxiomC{}
\RightLabel{$\delta_3$}
\DeduceC{$!B$}
\RightLabel{\Elim{\lif}}
\BinaryInfC{$!C$}
\DisplayProof
\]
To remove this cut
we replace the subproof~$\delta_1$ in~$\delta$ by the following:
\[
\AxiomC{}
\RightLabel{$\delta_3$}
\DeduceC{$!B$}
\RightLabel{$\Subst{\delta_2}{\delta_3}{!B^x}$}
\DeduceC{$!C$}
\DisplayProof
\]
By $\Subst{\delta_2}{\delta_3}{!B^x}$ we mean the !!{proof} that
results from~$\delta_2$ by replacing every assumption~$!B$ that's
labelled~$x$ with the entire !!{proof}~$\delta_3$. This !!{proof} no
longer has open assumptions of the form~$!B$ with label~$x$. It has no
other assumptions labelled~$x$, since we assume, as usual, that no two
!!{formula}s occur as assumptions with the same label. Hence, the new
!!{proof} has the same open assumptions as~$\delta_1$, every inference
in~$\delta_2$ that discharges assumptions discharges the same
assumptions in~$\Subst{\delta_2}{\delta_3}{!B^x}$, and every inference
in~$\delta_3$ that discharges assumptions results in possibly multiple
copies, which discharge the corresponding assumptions
in~$\Subst{\delta_2}{\delta_3}{!B^x}$.

Any cut that lies entirely within $\delta_2$, $\delta_3$, or
$\delta_4$, or entirely outside~$\delta_1$ in~$\delta$ is unchanged:
it still involves the same !!{formula}s, goes through the same
inferences, so in particular has the same rank and length as the
corresponding cut in~$\delta$. We have removed a maximal cut of
length~$1$ and so either have reduced the cut length, or reduced the
cut rank (if the cut length of~$\delta$ was~$1$ and this was the only
maximal cut).

However, two things may happen:
\begin{enumerate}
  \item By replacing all assumptions $!B^x$ by $\delta_3$, we have
  multiplied all the cuts in~$\delta_3$. If some of those were maximal
  cuts, the cut length of the new !!{proof} may be greater than the cut
  length of the old !!{proof}. We have to guarantee that this does not
  happen. We do so by applying reduction conversions only to
  \emph{rightmost} maximal cuts, that is, to cuts which are such that
  no maximal cut lies on a branch through~$\delta$ to the right of the
  cut being reduced. If there were a cut in~$\delta_3$ it would be to
  the right of the cut being reduced here, and so we wouldn't be
  working on a rightmost cut. By always selecting rightmost cuts, we
  guarantee that either we decrease the cut length of the !!{proof},
  or even reduce the cut rank.
  \item If an assumption $!B^x$ in $\delta_2$ is the major premise of
  !!a{elimination} inference, and the last inference of~$\delta_3$ is
  \Elim{\lor}, \Elim{\lexists}, or !!a{introduction} inference, we
  have generated a new cut by grafting~$\delta_3$ onto $\delta_2$ at
  such an assumption. What was just an assumption in~$\delta_2$ is now
  a cut of length~$1$ (the conclusion of !!a{introduction} inference
  and major premise of !!a{elimination} inference) or the last
  !!{formula} occurrence of a cut segment (major premise of
  !!a{elimination} inference and conclusion of~\Elim{\lor} or
  \Elim{\lexists}). If $B^x$ was a minor premise of \Elim{\lor} or
  \Elim{\lexists}, it already was the first !!{formula} occurrence of
  a segment, and possibly of a cut segment. In the new !!{proof} that
  segment may now be longer. However, observe that the !!{formula}
  making up this new cut or existing cut segment is~$!B$, and
  $\depth{!B} < \depth{!B \lif !C}$. Hence, although we have
  potentially added cuts or increased the length of cuts, none of
  these are maximal cuts. Thus we have either decreased the cut rank,
  or, if there are cuts of the same rank remaining, have decreased the
  cut length.
\end{enumerate}

The process of reducing cut of length~$1$ is called a \emph{reduction
conversion} (or \emph{detour conversion}). The
patterns for some of the rule pairs is given in
\olref{tab:reduction-conversions}. (The remaining are left as
exercises.) 

\begin{table}
\begin{multline*}
\bottomAlignProof
\AxiomC{$\Discharge{!B}{x}$}
\RightLabel{$\delta_2$}
\shortDeduce
\DeduceC{$!C$}
\DischargeRule{\Intro{\lif}}{x}
\UnaryInfC{$!B \lif !C$}
\AxiomC{}
\RightLabel{$\delta_3$}
\shortDeduce
\DeduceC{$!B$}
\RightLabel{\Elim{\lif}}
\BinaryInfC{$!C$}
\DisplayProof
\quad \longrightarrow\quad
\shoveright{\bottomAlignProof
\AxiomC{}
\RightLabel{$\delta_3$}
\shortDeduce
\DeduceC{$!B$}
\RightLabel{$\Subst{\delta_2}{\delta_3}{!B^x}$}
\shortDeduce
\DeduceC{$!C$}
\DisplayProof}
\\[4ex]
\shoveleft{\bottomAlignProof
\AxiomC{}
\RightLabel{$\delta_2$}
\shortDeduce
\DeduceC{$!B$}
\AxiomC{}
\RightLabel{$\delta_3$}
\shortDeduce
\DeduceC{$!C$}
\RightLabel{\Intro{\land}}
\BinaryInfC{$!B \land !C$}
\RightLabel{\Elim{\land}[1]}
\UnaryInfC{$!B$}
\DisplayProof
\quad \longrightarrow\quad
\bottomAlignProof
\AxiomC{}
\RightLabel{$\delta_2$}
\shortDeduce
\DeduceC{$!B$}
\DisplayProof}\qquad
\shoveright{\bottomAlignProof
\AxiomC{}
\RightLabel{$\delta_2$}
\shortDeduce
\DeduceC{$!B$}
\AxiomC{}
\RightLabel{$\delta_3$}
\shortDeduce
\DeduceC{$!C$}
\RightLabel{\Intro{\land}}
\BinaryInfC{$!B \land !C$}
\RightLabel{\Elim{\land}[2]}
\UnaryInfC{$!C$}
\DisplayProof
\quad \longrightarrow\quad
\bottomAlignProof
\AxiomC{}
\RightLabel{$\delta_3$}
\shortDeduce
\DeduceC{$!C$}
\DisplayProof}
\\[4ex]
\bottomAlignProof
\AxiomC{}
\RightLabel{$\delta_2$}
\shortDeduce
\DeduceC{$!B$}
\RightLabel{\Intro{\lor}[1]}
\UnaryInfC{$!B \lor !C$}
\AxiomC{$\Discharge{!B}{x}$}
\RightLabel{$\delta_3$}
\shortDeduce
\DeduceC{$!D$}
\AxiomC{$\Discharge{!C}{y}$}
\RightLabel{$\delta_4$}
\shortDeduce
\DeduceC{$!D$}
\DischargeRule{\Elim{\lor}}{x\,y}
\TrinaryInfC{$!D$}
\DisplayProof
\quad \longrightarrow\quad
\shoveright{\bottomAlignProof
\AxiomC{}
\RightLabel{$\delta_2$}
\shortDeduce
\DeduceC{$!B$}
\RightLabel{$\Subst{\delta_3}{\delta_2}{!B^x}$}
\shortDeduce
\DeduceC{$!D$}
\DisplayProof}
\\[2ex]
\shoveleft{\bottomAlignProof
\AxiomC{}
\RightLabel{$\delta_2$}
\shortDeduce
\DeduceC{$!B(c)$}
\RightLabel{\Intro{\lforall}}
\UnaryInfC{$\lforall[x][!B(x)]$}
\RightLabel{\Elim{\lforall}}
\UnaryInfC{$!B(t)$}
\DisplayProof
\quad \longrightarrow\quad
\bottomAlignProof
\AxiomC{}
\RightLabel{$\Subst{\delta_2}{t}{c}$}
\shortDeduce
\DeduceC{$!B(t)$}
\DisplayProof}
\\[4ex]
\shoveright{\bottomAlignProof
\AxiomC{}
\RightLabel{$\delta_2$}
\shortDeduce
\DeduceC{$!B(t)$}
\RightLabel{\Intro{\lexists}}
\UnaryInfC{$\lexists[x][!B(x)]$}
\AxiomC{$\Discharge{!B(c)}{x}$}
\RightLabel{$\delta_3$}
\shortDeduce
\DeduceC{$!D$}
\DischargeRule{\Elim{\lexists}}{x}
\BinaryInfC{$!D$}
\DisplayProof
\quad \longrightarrow\quad
\bottomAlignProof
\AxiomC{}
\RightLabel{$\delta_2$}
\shortDeduce
\DeduceC{$!B(t)$}
\RightLabel{$\Subst{\Subst{\delta_3}{t}{c}}{\delta_2}{!B(t)^x}$}
\shortDeduce
\DeduceC{$!D$}
\DisplayProof}
\end{multline*}
\caption{Some reduction conversions for \Log{N1i}.}
\ollabel{tab:reduction-conversions}
\end{table}

\begin{prob}
Give the reduction conversion for \Intro{\lnot} followed by
\Elim{\lnot}.
\end{prob}

There is one case still to be considered. The definition of cut
includes the case where the length of the cut is~$1$ and $!A$ is the
conclusion of~\FalseInt{} as well as the major premise of
!!a{elimination} inference. This case is dealt with in a similar way
as when $!A$ is the conclusion of !!a{introduction} rule. For
instance, replace
\[
\bottomAlignProof
\AxiomC{}
\RightLabel{$\delta_2$}
\DeduceC{$\lfalse$}
\UnaryInfC{$!B \lif !C$}
\AxiomC{}
\RightLabel{$\delta_3$}
\DeduceC{$!B$}
\RightLabel{\Elim{\lif}}
\BinaryInfC{$!C$}
\DisplayProof
\quad\text{by}\quad
\bottomAlignProof
\AxiomC{}
\RightLabel{$\delta_2$}
\DeduceC{$\lfalse$}
\UnaryInfC{$!C$}
\DisplayProof
\]
We must be a little careful in the cases where $!A \ident (!B \lor
!C)$ or $!A \ident \lexists[x][!B(x)]$. E.g., if we replaced
\[
\bottomAlignProof
\AxiomC{}
\RightLabel{$\delta_2$}
\DeduceC{$\lfalse$}
\RightLabel{\FalseInt}
\UnaryInfC{$!B \lor !C$}
\AxiomC{$\Discharge{!B}{x}$}
\RightLabel{$\delta_3$}
\DeduceC{$!D$}
\AxiomC{$\Discharge{!C}{x}$}
\RightLabel{$\delta_4$}
\DeduceC{$!D$}
\RightLabel{\Elim{\lor}}
\TrinaryInfC{$!D$}
\DisplayProof
\quad\text{by}\quad
\bottomAlignProof
\AxiomC{}
\RightLabel{$\delta_2$}
\DeduceC{$\lfalse$}
\RightLabel{\FalseInt}
\UnaryInfC{$!D$}
\DisplayProof
\]
we might be introducing a new cut with cut !!{formula}~$!D$, and there
is no guarantee that $\depth{!D} < \depth{!B \lor !C}$!{} So we must
proceed the same way here as we did for an \Intro{\lor}/\Elim{\lor}
reduction and replace with
\[
\bottomAlignProof
\AxiomC{}
\RightLabel{$\delta_2$}
\DeduceC{$\lfalse$}
\RightLabel{\FalseInt}
\UnaryInfC{$!B$}
\RightLabel{$\delta_3$}
\DeduceC{$!D$}
\DisplayProof
\quad\text{or}\quad
\bottomAlignProof
\AxiomC{}
\RightLabel{$\delta_2$}
\DeduceC{$\lfalse$}
\RightLabel{\FalseInt}
\UnaryInfC{$!C$}
\RightLabel{$\delta_4$}
\DeduceC{$!D$}
\DisplayProof
\]

\end{document}
